% Historical source (removed by author request): references/_OLD_Übungsaufgaben/Übungsblatt 2 - Lineare Gleichungssysteme/Blatt 2.tex:96
\begin{ExerciseBlock}
\ExerciseHeading{02}{Lösungsmenge mit Gauß-Verfahren}{linear_systems_dependent_equations}
Überprüfen Sie mit dem Gauß-Verfahren die Lösungsmenge eines LGS, bei dem

\begin{enumerate}[label=\alph*)]
\item eine Gleichung die Summe zweier anderer Gleichungen ist.
	\begin{equation*}
	\begin{split}
		a x_1 + b x_2 + c x_3 &= d \\
 		e x_1 + f x_2 + g x_3 &= h  \\
 		(a+e) x_1 + (b+f) x_2 + (c+g) x_3 &= d+h
	\end{split}
	\end{equation*}
\item eine Gleichung ein Vielfaches einer anderen Gleichung ist.
	\begin{equation*}
	\begin{split}
		a x_1 + b x_2 + c x_3 &= d \\
 		e x_1 + f x_2 + g x_3 &= h  \\
 		\lambda a x_1 + \lambda b x_2 +  \lambda c x_3 &= \lambda d
	\end{split}
	\end{equation*}
	Vergleichen Sie diese Erkenntnisse mit Aufgabe 1.
\end{enumerate}

\end{ExerciseBlock}
\UseExerciseSolution{linear_systems_dependent_equations}
